\subsection{Path-connected spaces}

DEFINITION 3. --- A topological space X is said to be path-connected if for every pair \((x,y)\) of points of X, there exists a path with origin x and endpoint y.

Example 1. --- Every interval of R, every convex subset of a numerical space or, more generally, of a topological vector space over R is path-connected.

PROPOSITION 3. --- The image under a continuous map of a path-connected space is path-connected.

Let X be a path-connected space, \(f: X \to Y\) a continuous map. Let x and y be two points of \(f(X)\); let \(x'\) and \(y'\) be points of X such that \(f(x') = x\) and \(f(y') = y\). There exists a path c in X whose origin is \(x'\) and endpoint \(y'\). Then, \(f \circ c\) is a path in \(f(X)\) with origin x and endpoint y.

PROPOSITION 4. --- A path-connected topological space is connected.

Since every interval of R is connected, the image of a path is connected (TG, I, p. 82, prop. 4). The empty space is connected. Since a nonempty path-connected space is the union of the images of paths issuing from one of its points, it is connected (TG, I, p. 81, prop. 2).

In view of the identification (III, p. 256, remark 1) of paths in X with homotopies between maps of a topological space P reduced to a point into X, proposition 1 of III, p. 230 implies:

PROPOSITION 5. --- In a topological space, the relation “there exists a path with origin x and endpoint y” is an equivalence relation.

The path-connected components of a topological space X are called the equivalence classes for the above relation. A path-connected component is a path-connected space. Let x be a point of X. The path-connected component of x is the union of the path-connected subspaces of X containing x. It is also the union of the images of paths with origin x in X.

Examples. --- 2) The union of a family of path-connected sets, whose intersection is nonempty, is path-connected (cf. TG, I, p. 81, prop. 2).

\begin{enumerate}
\def\labelenumi{\arabic{enumi})}
\setcounter{enumi}{2}
\tightlist
\item
  A subset of a numerical space (or, more generally, of a topological vector space over R) which is star-shaped at one of its points is path-connected.
\end{enumerate}

Let X be a topological space. We denote by \(\pi_{0}(X)\) the set of path-connected components of X. Let P be a topological space reduced to a point. The map from X to {[}P; X{]} which, to every point x of X, associates the homotopy class \(\varphi_{x}\) of the map \(f_{x}: P \to X\) with image x, defines by passage to the quotient a bijection, called canonical, from \(\pi_{0}(X)\) onto {[}P; X{]}. We have \(\pi_{0}(\varnothing) = \varnothing\) . For a nonempty topological space to be path-connected, it is necessary and sufficient that \(\pi_{0}(X)\) be a set with one element.

Let X and Y be topological spaces and \(f\colon\mathrm{X}\to\mathrm{Y}\) a continuous map. The image \(f(\mathrm{C})\) of every path-connected component C of X is path-connected (III, p. 258, prop. 3), hence contained in a path-connected component of Y. We denote by \(\pi_{0}(f)\colon\pi_{0}(\mathrm{X})\to\pi_{0}(\mathrm{Y})\) the map which, to a path-connected component C of X, associates the unique path-connected component C' of Y such that \(f(\mathrm{C})\subset\mathrm{C}'\). If one identifies \(\pi_0(\mathrm{X})\) and \(\pi_0(\mathrm{Y})\) with {[}P; X{]} and {[}P; Y{]} respectively, \(\pi_0(f)\) is identified with the map \(\chi \mapsto [f] \circ \chi\) from {[}P; X{]} to {[}P; Y{]}. In particular, if \(f\) and \(g\) are homotopic maps from X to Y, we have \(\pi_0(f) = \pi_0(g)\) (III, p. 230, prop. 2).

Let Z be a topological space and let \(g: Y \to Z\) be a continuous map; we have \(\pi_{0}(g \circ f) = \pi_{0}(g) \circ \pi_{0}(f)\) . In particular, if Z = X and if f and g are homotopy inverses of each other, we have \(\pi_{0}(g) \circ \pi_{0}(f) = \pi_{0}(\mathrm{Id}_{\mathrm{X}}) = \mathrm{Id}_{\pi_{0}(\mathrm{X})}\) and \(\pi_{0}(f) \circ \pi_{0}(g) = \pi_{0}(\mathrm{Id}_{\mathrm{Y}}) = \mathrm{Id}_{\pi_{0}(\mathrm{Y})}\) , which proves that \(\pi_{0}(f)\) and \(\pi_{0}(g)\) are inverse bijections of each other. A space homotopy equivalent to a path-connected space is therefore itself path-connected. In particular, a space homotopy equivalent to a point is path-connected.

PROPOSITION 6. --- Let \((\mathrm{Y}_{j})_{j\in\mathrm{J}}\) be a family of topological spaces. The map

\[
(\pi_ {0} (\mathrm{pr} _ {j})) \colon \pi_ {0} (\prod_ {j \in \mathrm{J}} \mathrm{Y} _ {j}) \to \prod_ {j \in \mathrm{J}} \pi_ {0} (\mathrm{Y} _ {j})
\]

is bijective. In particular, the product space of a family of path-connected spaces is path-connected.

This follows from Proposition 3 of III, p. 231, where one takes for X a space reduced to a single point.

